% 07-full-map.tex -- the whole map.  Formula F18, certificate V15.
% Input by mostlyharmless.tex.  No preamble here.

\section{The whole map}
\label{sec:full-map}

Everything a map needs is now in place, and assembling it is arithmetic: the sum over the
stored half of a map's coefficients is \eqref{eq:pairing-visibility-parity}, derived in
section~\ref{sec:parity} and not restated here, with the mode visibilities from
\eqref{eq:mode-visibility} --- the odd family's from \eqref{eq:odd-mode}, the even family's
from \eqref{eq:even-radial}.

Two properties are worth stating, because they are what an observer can check.

The first is Hermitian symmetry.  For a real map
\begin{equation}
  V(-u,-v) = \overline{V(u,v)} ,
  \label{eq:hermitian}
\end{equation}
so the real part of a visibility map is even and the imaginary part odd.  That is
visible by eye in \code{make mapgui-levels}, and it is checked directly.

The second is a bound,
\begin{equation}
  |V(u,v)| \le V(0,0) \qquad \text{when } I \ge 0 ,
  \label{eq:visibility-bound}
\end{equation}
since $|V| \le \int|I| = \int I = V(0,0)$ for a non-negative map.  The condition
matters: a truncated expansion of a discontinuous map dips below zero near its
edges, which is the truncation ringing of section~\ref{sec:full-map}, and there the
bound fails.
Stating it with its condition is the honest form.

\subsection{What was checked}
\label{sec:full-map-register}

\begin{center}
\begin{tabular}{lll}
\hline
tag & claim & rung \\
\hline
\vtag{15} & \eqref{eq:pairing-visibility-parity} against a 2-D FFT of the rendered image & numerical, low-frequency band \\
\hline
\end{tabular}
\end{center}

Certificate \vtag{15} renders a rotated starspot, transforms the image with a 2-D
FFT using the same planar measure, and compares.  It is the end-to-end anchor: the
coefficients, both parity families, the rotation and the measure all enter, and it
is the only check that would catch an error common to both families.  The comparison
is restricted to low frequencies, because the rim of the disk is a discontinuity and
its Gibbs ringing contaminates high frequencies at any resolution.  Hermitian
symmetry \eqref{eq:hermitian} is checked separately and much more tightly, being a
statement about the transform rather than about a resolution.

That ringing is the practical cost of the whole scheme, and it is worth stating as a rate
rather than as a picture.  A map whose brightness jumps --- a sharp-edged spot, a polar cap
--- has coefficients falling as a power of $l$, so it converges \emph{algebraically} and not
exponentially.  The error is then set by the Gibbs ring rather than by the smooth part: the
overshoot at the edge is a fixed fraction of the jump, only its width shrinks, so the mean
square falls as $l_{\max}^{-1}$ and the root mean square as $l_{\max}^{-1/2}$ --- to
roughly two thirds when the order doubles, which is what certificate \vtag{28} measures.  A
smooth map converges far faster, which is why the limb darkening of section~\ref{sec:models}
buys accuracy
as well as realism.  That rate is what an expansion in spherical harmonics costs, and it is
the reason the mode count is a knob rather than a formality.

\begin{itemize}
  \item[$\checkmark$] \code{check_convergence_rate} compares each truncation with $l = 64$
    on the rendered disc and reports the worst ratio between successive errors.  It is
    measured on the \emph{map} rather than on the visibility for a reason worth recording:
    every even-family mode of a visibility costs a Zernike expansion in exact rationals, so
    the cost grows as roughly the fourth power of the order, and a reference at high $l$ is
    unaffordable.  Both sides are in the same units, so the normalisation cancels in the
    ratio and only the rate survives --- which is the claim.
    \\[1pt] \vtag{28}\quad \emph{Numerical}, $l \le 32$
\end{itemize}
